\section{Conclusion and Future Work}
\label{sec:conclusion}

In this paper, we have critically examined the recently popularized paradigm of online test-time adaptation (TTA) for weight-space model merging. From the perspective of \textbf{The Methodologist}, we exposed two severe methodological weaknesses: the ``no-data'' strawman, which compares complex online TTA against an unoptimized uniform baseline while ignoring the practical availability of small validation sets, and the catastrophic fragility of unsupervised online objectives under non-i.i.d.\ and noisy target streams.

To address these limitations, we proposed \textbf{Offline Few-Shot Validation Tuning (OFS-Tune)}. By optimizing highly constrained, low-dimensional coefficient profiles offline on as few as 5 to 10 validation samples per task, OFS-Tune completely bypasses the need for test-time adaptation. Across a rigorous multi-seed evaluation, OFS-Tune consistently outperformed SOTA online TTA methods under standard streams, whilst demonstrating absolute, unwavering robustness under realistic adversarial distribution shifts. Furthermore, we validated the \emph{Overfitting-Optimizer Paradox}, showing that low-dimensional parameterizations (such as low-degree polynomials) act as vital low-pass filters that prevent validation overfitting in the scarce few-shot regime.

Our findings serve as a vital methodological course correction for the model-merging community. We strongly urge future research to move away from overly complex, fragile, and compute-heavy online adaptation schemes that chase minor performance gains under artificial, clean conditions. Instead, we call on researchers to evaluate their proposed methods against strong, simple baselines like few-shot validation tuning, and rigorously stress-test their models under adversarial deployment streams.

\subsection{Limitations and Future Work}
\label{sec:conclusion_limitations}
While our continuous simulation framework is carefully calibrated on empirical Vision Transformer (ViT-B/32) classification statistics, it remains a mathematical abstraction of deep neural landscapes. Real weight-merging environments exhibit complex, discontinuous topologies and dynamic batch characteristics. Future work should empirically validate OFS-Tune on physical deep networks and scale our low-dimensional parameterizations (such as low-rank matrices) to LLMs and image diffusion architectures. A detailed, rigorous discussion of our methodological limitations, simulation abstractions, and future research avenues is presented in Section \ref{appendix:limitations} of the Appendix.
